\section{Discussion}
\label{sec:discussion}

\subsection{Key Empirical and Analytical Insights}
Our simulation results align closely with the mathematical models, establishing three primary takeaways. First, yielding fundamentally relocates the performance bottleneck from compute cores to the network fabric. Because busy-waiting pins cores for the entire duration of a remote fetch, 8 worker cores cap throughput at 1.77~MRPS even though the 100~GbE link can carry substantially more. By freeing cores to continue processing local requests and issuing subsequent fetches, yielding drives system throughput to the physical 2.50~MRPS link bound. Second, yielding prevents early queue buildup below saturation: at 1.5~MRPS, yielding cuts P99.9 tail latency roughly in half compared to spinning. Third, the advantage relies strictly on microsecond-scale context switching; the inequality $2t_s < t_f$ serves as a practical litmus test for whether yielding is worth implementing on a given system.

From a resource provisioning perspective, busy-waiting wastes over half (54\%) of active core cycles spinning idly on RDMA completions at saturation. Matching the 2.50~MRPS throughput of an 8-worker yielding node would require approximately 11 spinning cores, representing a substantial hardware overhead.

\subsection{Pathologies of Pure Yielding and the Hybrid Solution}
Our sensitivity sweeps demonstrate that pure yielding is not universally optimal:
\begin{enumerate}
\item \emph{High Context-Switch Latency}: When thread switching exceeds $t_s \approx 1.5~\mu$s, context-switch overhead outweighs the fetch duration, driving yielding throughput below spinning.
\item \emph{Light Workload Penalties}: At low arrival rates, cores are largely idle. Yielding on every miss incurs two context switches ($2 \times 0.02 = 0.04~\mu$s) without providing any concurrency benefit, adding an unnecessary 20~ns latency overhead. In systems with slower thread packages, this overhead expands to several microseconds.
\item \emph{High Cache Hit Rates}: When local hit rates reach 80\%, page faults become infrequent, causing spinning and yielding to converge to identical performance ceilings.
\end{enumerate}

Our \emph{Load-Aware Adaptive Hybrid Yielding} policy eliminates these drawbacks. By checking queue occupancy locally before deciding to yield, workers spin when the server is uncongested and switch only when backlog accumulates. This provides the low latency of spinning at light load, the high saturation throughput of yielding under stress, and complete immunity to context-switch degradation up to $t_s = 1.5~\mu$s.

\subsection{Explaining Divergences from Published Adios Data}
Table~\ref{tab:paperlink} details the points of comparison between our simulation and the published Adios paper. While both agree on qualitative trends, numerical values diverge in three specific areas:
\begin{itemize}
\item \emph{Busy-wait saturation}: Adios reported $\approx 1.4$~MRPS on its microbenchmark, whereas our simulator yields 1.77~MRPS. This is expected because real-world NIC doorbell ringing and lock contention add overheads not captured in our two fixed CPU service times.
\item \emph{PF-aware dispatch gain}: Adios achieved up to 7.5\% (Memcached) and 27\% (RocksDB) tail latency improvements from congestion-aware dispatching, whereas our model observes $\le 1.2\%$. Real RDMA fabrics experience bursty arrivals and variable service times where queue imbalance is pronounced; in our deterministic-time model, the shared link remains the dominant bottleneck, leaving little per-QP delay for the dispatcher to eliminate.
\item \emph{Memory pressure degradation}: Real OS paging incurs lock contention, frame allocation stalls, and synchronous writebacks under memory pressure, which explains why Adios's measured throughput dropped faster when DRAM shrank.
\end{itemize}

\subsection{Assumptions and Threats to Validity}
To maintain focus on core scheduling dynamics, our model makes several simplifying assumptions:
\begin{itemize}
\item \textbf{Deterministic service latencies}: CPU and network delays are modeled as fixed constants rather than stochastic distributions, which likely compresses the tail compared to real systems.
\item \textbf{Idealized dispatcher}: The central dispatcher runs with zero CPU overhead and infinite throughput, abstracting away central lock contention that would limit scalability beyond 10--12 cores.
\item \textbf{Absence of reclaimer and prefetcher}: Free page frames are assumed to be immediately available; we do not model page eviction delays or multi-page prefetching pipelines.
\item \textbf{Unmerged page requests}: Concurrent misses to the same page produce duplicate RDMA fetches rather than merging into a single transfer.
\item \textbf{Finite in-flight pool}: We cap concurrent in-flight requests at 4,096 to prevent memory leaks during overload, which is smaller than Adios's 131,072 pool.
\end{itemize}

\subsection{Methodological Reflections}
During development, comparing simulation plateaus against closed-form bounds caught two significant implementation bugs: an arrival-timestamp windowing error that distorted saturation throughput, and an unbounded request pool that allowed physically impossible link throughputs ($>2.50$~MRPS). Deriving simple capacity equations proved essential for verifying simulation correctness before interpreting graphs.

\subsection{Directions for Future Research}
Several extensions would build naturally on this study. Simulating variable core counts ($N = 2$--16) and remote fetch times ($t_f = 2$--6~$\mu$s) would directly validate our analytical predictions. Incorporating heavy-tailed request service times and stochastic network jitter would better capture tail effects. Finally, evaluating dispatcher CPU overhead and implementing a multi-queue decentralized dispatcher would test scalability limits for modern 64-core and 128-core servers.
